\subsection{原函数的定义}

向量函数 \(\mathbf{f}\)（定义在区间 \(\mathrm{I} \subset \mathbf{R}\) 上），除非满足相当苛刻的条件，否则不可能是一个向量函数 \(\mathbf{g}\)（定义且连续在 \(\mathrm{I}\) 上）在此区间的每一点处的导数：例如，若 \(\mathbf{f}\) 在一点 \(x_0\)（内于 \(\mathrm{I}\)）处具有右极限与左极限，则 \(\mathbf{f}\) 必须在点 \(x_0\) 处连续，这由 \(\mathrm{I}\), p.~18 的命题 6 得出；由此可知，若取区间 \([-1, 1]\) 作为 \(\mathrm{I}\)，并取 \(f\) 为等于 \(-1\)（在 \([-1, 0[\) 上）、等于 \(+1\)（在 \([0, 1]\) 上）的实函数，则 \(f\) 不是 \(\mathrm{I}\) 上任何连续函数的导数；尽管如此，函数 \(|x|\) 都以 \(f(x)\) 为导数（在每个点 \(\neq 0\) 处）；于是人们被引导作出以下定义：

定义 1. 给定定义在区间 I ⊂ R 上的向量函数 f，我们称定义在 I 上的函数 g 是 f 的原函数，如果 g 在 I 上连续，并且在 I 的某个可数子集（相对于 I）的余集的每个点 x 处具有等于 f(x) 的导数。

若此外 \(\mathbf{g}\) 还具有等于 \(\mathbf{f}(x)\) 的导数（在 I 的每个点 \(x\) 处），则称 \(\mathbf{g}\) 是 \(\mathbf{f}\) 的严格原函数。

有了这个定义，就可看出上面考虑的实函数 \(f\) 具有一个等于 \(|x|\) 的原函数。

显然，若 f 在 I 上具有原函数，则 f 的每个原函数也是任何除在 I 的某个可数子集的点处外都等于 f 的函数的原函数。在用语从宽的意义下，人们也说到只定义在 I 的某个可数子集（相对于 I）的余集上的函数 \(f_{0}\) 在 I 上的原函数：它将是每个定义在 I 上且等于 \(f_{0}\)（在 \(f_{0}\) 有定义的点处）的函数 f 的原函数。

命题 1. 设 f 是定义在 I 上、取值于 E 中的向量函数；若 f 在 I 上具有原函数 g，则 f 在 I 上的原函数所成之集等同于诸函数 \(g + a\) 所成之集，其中 a 是取值于 E 中的常值函数。

事实上，显然 \(\mathbf{g} + \mathbf{a}\) 都是 \(\mathbf{f}\) 的原函数（对任何 \(\mathbf{a} \in E\)）；另一方面，若 \(\mathbf{g}_1\) 是 \(\mathbf{f}\) 的原函数，则 \(\mathbf{g}_1 - \mathbf{g}\) 除在 I 的某个可数子集的点处外具有等于 0 的导数，从而是常数（I, p.~17, 推论）。

我们称函数 f 的原函数（当它们存在时）是“相差一个加性常数”确定的。为无歧义地定义 f 的一个原函数，只需（任意地）给它指定在一点 \(x_{0} \in I\) 处的值；特别地，存在唯一的一个原函数 g，使得 \(\mathbf{g}(x_{0}) = 0\)；对 f 的每个原函数 h 有 \(\mathbf{g}(x) = \mathbf{h}(x) - \mathbf{h}(x_{0})\)。

